\section{Central nefness and exact FS functions of power-tensor perturbations}
\label{sec:native-tc-nef-fs-tensor-lines}

Fix a test configuration $\pi\colon\mathcal X\to\mathbf A^1$ of $X$, and let
$\mathcal X_0=\mathcal X\times_{\mathbf A^1}\{0\}$ be its whole scheme-theoretic central
fibre, with inclusion $i_0$.  Let $(M,\alpha,e)$, $(H,\beta,f)$ be linearized lines with
generic lines $L,K$ (Lemma~\ref{lem:native-tc-nef-fs-whole-tensor-comparisons}).

\begin{lemma}[The central fibre is projective {\cite[\S1.2]{BJ22}}]
\label{lem:native-tc-nef-fs-central-projective}
\leanok
\lean{EnvelopeProperty.NativeTestConfiguration.TestConfiguration.centralProjection_isProjective}
\uses{lem:intrinsic-original-normalized-central-native-projective}
The structure map $\mathcal X_0\to\operatorname{Spec}k$ is projective.
\end{lemma}

\begin{proof}
The projection $\pi$ factors as a closed immersion into a projective bundle over
$\mathbf A^1$ followed by the bundle projection.  Base change along
$0\colon\operatorname{Spec}k\to\mathbf A^1$ gives such a factorization of
$\mathcal X_0\to\operatorname{Spec}k$, which is therefore projective.
\end{proof}

\begin{theorem}[Central ampleness of the perturbation {\cite[Lemma~5.5]{FHX}}]
\label{thm:native-tc-nef-fs-central-tensor-ample}
\leanok
\lean{EnvelopeProperty.NativeTestConfiguration.TestConfiguration.central_ample_power_tensor_of_nef}
\uses{lem:native-tc-nef-fs-central-projective, thm:native-nef-plus-ample}
Suppose $i_0^*M$ is nef, i.e.\ has nonnegative degree on every closed integral curve of
$\mathcal X_0$, and $i_0^*H$ is ample.  Then
$i_0^*(M^{\otimes b}\otimes H^{\otimes a})$ is ample for every natural $b$ and every
positive natural $a$.
\end{theorem}

\begin{proof}
By Lemma~\ref{lem:native-tc-nef-fs-central-projective} the central fibre is projective over
$k$.  The class of $(i_0^*M)^{\otimes b}$ is $b$ times the class of $i_0^*M$, so it is
nef.  A positive power of an ample line on a projective scheme is ample: its pullback to
the closure of every point is a positive power of an ample line on an integral projective
scheme, and ampleness may be tested on these closures.  The tensor product of a nef line
and an ample line on a projective $k$-scheme is ample (Theorem~\ref{thm:native-nef-plus-ample});
this applies to $\mathcal X_0$ even when it is reducible or nonreduced.  Finally the
tensorator and the power comparison of $i_0^*$ identify
$(i_0^*M)^{\otimes b}\otimes(i_0^*H)^{\otimes a}$ with
$i_0^*(M^{\otimes b}\otimes H^{\otimes a})$, and ampleness is invariant under
isomorphism.
\end{proof}

\begin{theorem}[Exact FS realization of the perturbed potential {\cite[Lemma~5.5]{FHX}}]
\label{thm:native-tc-nef-fs-perturbed-model-exact}
\leanok
\lean{EnvelopeProperty.NativeTestConfiguration.TestConfiguration.isFinFS_power_tensor_modelPotential_of_central_nef}
\uses{thm:native-tc-nef-fs-central-tensor-ample, thm:central-ample-original-exact-finite-fs,
lem:native-tc-nef-fs-whole-tensor-comparisons}
Assume $k$ is infinite and $\mathcal X$ is integral and dominant over $\mathbf A^1$.  Under
the hypotheses of Theorem~\ref{thm:native-tc-nef-fs-central-tensor-ample}, with $a>0$, the native
model potential of $M^{\otimes b}\otimes H^{\otimes a}$, for the linearization and
fibre-one identification of Lemma~\ref{lem:native-tc-nef-fs-whole-tensor-comparisons}, is a
finite rational Fubini--Study function of class $1$ for
$L^{\otimes b}\otimes K^{\otimes a}$.
\end{theorem}

\begin{proof}
By Theorem~\ref{thm:native-tc-nef-fs-central-tensor-ample} the restriction of
$M^{\otimes b}\otimes H^{\otimes a}$ to the central fibre is ample.  Apply
Theorem~\ref{thm:central-ample-original-exact-finite-fs} to this linearized line: a positive
power is generated by finitely many global sections, so the base-ideal correction
vanishes and the weighted FS function of these generators equals the native model
potential at every point.
\end{proof}

\begin{remark}
The identification of this potential with $b\varphi_M+a\varphi_H$, where
$\varphi_M,\varphi_H$ are the native model potentials of $M$ and $H$, is
Theorem~\ref{thm:native-tc-nef-fs-power-tensor-native-potential}.  The construction of an
auxiliary $H$ with $i_0^*H$ ample from the PL model is not part of this chapter.
\end{remark}
